\chapter{Perturbative Algebraic QFT On Curved Backgrounds}
\label{ch:curved-perturbative-algebraic-qft}

\ConventionsMixed

This chapter develops perturbative algebraic QFT on curved backgrounds as a
local construction of interacting fields from Hadamard two-point functions,
microcausal functionals, time-ordered products, and causal factorization.
The framework is perturbative, but its locality and covariance conditions are
nonperturbative constraints on the allowed coefficients.

\section{Microcausal Functionals}

Let \(M\) be globally hyperbolic and let \(\mathcal E=C^\infty(M)\) be the
space of scalar field configurations.  A functional \(F:\mathcal E\to\C\) is
microcausal when its functional derivatives
\[
  F^{(n)}(\phi)\in\mathcal E'(M^n)
\]
have wavefront sets avoiding products of closed future or past light cones:
\[
  \operatorname{WF}(F^{(n)}(\phi))
  \cap
  \left(\overline V_+^{\,n}\cup\overline V_-^{\,n}\right)
  =
  \varnothing .
\]
This condition makes the contractions with Hadamard distributions defined by
the wavefront-set product criterion.

\section{Star Product}

Choose a Hadamard two-point function \(H\).  The star product of microcausal
functionals is
\[
  F\star_H G
  =
  \sum_{n=0}^{\infty}
  \frac{\hbar^n}{n!}
  \left\langle
    F^{(n)},H^{\otimes n}G^{(n)}
  \right\rangle .
\]
For polynomial functionals the sum is finite.  Associativity follows from the
fact that the bidifferential contraction operator
\[
  \Gamma_H=\left\langle H,\frac{\delta}{\delta\phi}\otimes
  \frac{\delta}{\delta\phi}\right\rangle
\]
has commuting copies on tensor products, so
\[
  \exp(\hbar\Gamma_H)
\]
defines an associative deformation of pointwise multiplication.

\section{Change Of Hadamard Function}

If \(H'\) is another Hadamard two-point function, then
\[
  w=H'-H
\]
is smooth.  Define
\[
  \alpha_w
  =
  \exp\!\left(
    \frac{\hbar}{2}
    \left\langle w,\frac{\delta^2}{\delta\phi^2}\right\rangle
  \right).
\]
Then
\[
  \alpha_w(F\star_H G)=\alpha_w(F)\star_{H'}\alpha_w(G).
\]
Thus the abstract algebra is independent of the particular Hadamard
parametrix up to a canonical isomorphism determined by the smooth difference.

\section{Time-Ordered Products}

Let \(T_n\) assign distributions on \(M^n\) to \(n\) local functionals.  Away
from coincident diagonals, \(T_n\) is determined by causal ordering and the
Feynman propagator.  The construction problem is to extend \(T_n\) to the
diagonals while preserving:
\[
\begin{gathered}
  \text{local covariance},\qquad
  \text{causal factorization},\\
  \text{field independence},\\
  \text{scaling degree bounds},\qquad
  \text{unitarity and Ward identities}.
\end{gathered}
\]
The extension freedom is a finite sum of local curvature and coupling
polynomials multiplying derivatives of delta distributions on diagonals.

\section{Causal Factorization}

If the supports of \(F_1,\ldots,F_k\) lie nowhere in the past of the supports
of \(F_{k+1},\ldots,F_n\), causal factorization states
\[
  T_n(F_1,\ldots,F_n)
  =
  T_k(F_1,\ldots,F_k)\star
  T_{n-k}(F_{k+1},\ldots,F_n).
\]
This identity determines \(T_n\) outside the total diagonal from lower orders.
Renormalization is precisely the choice of extension to the total and partial
diagonals compatible with the listed constraints.

\section{Interacting Fields}

Let \(V\) be an interaction functional with compact time support.  The
relative S-matrix is
\[
  S_V(F)=S(V)^{-1}\star S(V+F),
\]
where
\[
  S(F)=T\exp\!\left(\frac{\ii}{\hbar}F\right).
\]
The interacting field associated to a local functional \(F\) is the
Bogoliubov derivative
\[
  R_V(F)
  =
  \left.\frac{\hbar}{\ii}\frac{d}{d\lambda}
  S_V(\lambda F)\right|_{\lambda=0}.
\]
Changing the time cutoff of \(V\) outside the causal domain of \(F\) gives
isomorphic interacting algebras by causal factorization.

\section{Renormalization Group}

Two choices of time-ordered products \(T\) and \(\widetilde T\) are related
by maps \(Z\) on local functionals:
\[
  \widetilde S(F)=S(Z(F)).
\]
The map \(Z\) is local, begins with \(Z(F)=F+O(F^2)\), preserves the constant
term, and has coefficients supported on coincident points.  In curved
backgrounds these coefficients include curvature terms.  This is the
perturbative algebraic form of local counterterm freedom.

\begin{openproblem}[Nonperturbative curved-background comparison]
\label{op:nonperturbative-paqft-comparison}
Relate perturbative algebraic QFT on curved backgrounds to nonperturbative
constructions by comparing local algebras, Hadamard states, stress tensors,
renormalized composite fields, anomaly terms, and scaling limits in examples
where both sides are available.
\end{openproblem}
